% ==========================================================================
% TEMPLATE DE PROPOSTA DE TRABALHO - ESTAGIO EM DOCENCIA (PPGCC/DCC/UFMG)
% Estrutura exigida pelas Diretrizes (secao "Proposta de trabalho"):
%   1. Identificacao   2. Sobre a Disciplina   3. Atividades   4. Cronograma
%
% COMO USAR:
%  1) Preencha os dados nos comandos \newcommand logo abaixo.
%  2) Leia as caixas azuis "INSTRUCAO" e substitua os textos entre
%     colchetes [ ... ] pelo seu conteudo.
%  3) Ao terminar, troque \mostrarinstrucoestrue por \mostrarinstrucoesfalse.
%  4) Compile: pdflatex main (duas vezes, por causa da longtable)
%  5) Envie ao Moodle em PDF (nome: NomeCompletoDoAluno.pdf) e avise por
%     e-mail o coordenador da disciplina, com o professor em copia.
% ==========================================================================
\documentclass[12pt, a4paper]{article}
\usepackage[utf8]{inputenc}
\usepackage[T1]{fontenc}
\usepackage[brazilian]{babel}
\usepackage{longtable}
\usepackage{amsmath}
\usepackage{xcolor}
\usepackage[most]{tcolorbox}
\usepackage{hyperref}
\usepackage{indentfirst}
\usepackage{enumitem}
\usepackage{titling}
\setlength{\droptitle}{-10em}
\usepackage{geometry}
\geometry{a4paper, margin=2.5cm}

\definecolor{dccblue}{RGB}{20, 45, 110}
\definecolor{instrbg}{RGB}{232, 240, 254}

% ---------------------------- DADOS GERAIS ---------------------------------
\newcommand{\NomeEstagiario}{[Nome completo do Estagiário]}
\newcommand{\NivelEstagiario}{[Mestrando/Doutorando] em Ciência da Computação}
\newcommand{\EstagioNivel}{[Estágio em Docência I ou II]}
\newcommand{\NomeDisciplina}{[Nome da Disciplina]}
\newcommand{\CodigoDisciplina}{[CODIGO]}
\newcommand{\Turma}{[Turma]}
\newcommand{\Horarios}{[dias da semana e horários das aulas]}
\newcommand{\Semestre}{[AAAA/S]}
\newcommand{\NomeProfessor}{[Nome completo do Professor Responsável]}
\newcommand{\TituloProfessor}{[Titulação, ex.: Ph.D. em Ciência da Computação]}

% ------------------ CAIXAS DE INSTRUCAO (ocultaveis) -----------------------
\newif\ifmostrarinstrucoes
\mostrarinstrucoestrue   % <- troque por \mostrarinstrucoesfalse na versao final
\newcommand{\instrucao}[1]{%
  \ifmostrarinstrucoes
  \begin{tcolorbox}[colback=instrbg, colframe=dccblue, arc=2mm, boxrule=0.6pt,
    title=\textbf{INSTRUÇÃO (apagar na versão final)}, fonttitle=\small]
  \small #1
  \end{tcolorbox}
  \fi}

\title{\textbf{\EstagioNivel{} \\ Proposta de Trabalho}}
\author{\NomeEstagiario}
\date{[Data de submissão]}

\begin{document}
\maketitle

% ============================ 1. IDENTIFICACAO ==============================
\section{Identificação} \label{sec:identificacao}
\instrucao{Item 1 das Diretrizes. A proposta deve ser previamente acertada com o professor responsável pela turma à qual você foi alocado. Confira os dados na planilha oficial de alocação do semestre.}

\textbf{Título:} \EstagioNivel{} --- \NomeDisciplina

\textbf{Disciplina:} \NomeDisciplina

\textbf{Código:} \CodigoDisciplina

\textbf{Horários:} \Horarios

\textbf{Professor(a) Responsável:} \NomeProfessor

\textbf{Turma:} \Turma

\textbf{Período:} \Semestre

% ========================== 2. SOBRE A DISCIPLINA ===========================
\section{Sobre a Disciplina} \label{sec:disciplina}
\instrucao{Item 2 das Diretrizes: descreva a disciplina conforme o plano de ensino oficial e a condução adotada pelo professor. Consulte o plano de ensino e o professor para garantir fidelidade.}

\subsection{Condução pelo professor}
\instrucao{1 ou 2 parágrafos: modalidade (teórica, prática, teórico-prática), número de aulas semanais, uso do Moodle, laboratórios, trabalhos práticos etc.}
[Descrição da forma de condução da disciplina.]

\subsection{Ementa e Programa da Disciplina}
\instrucao{Copie a ementa oficial e o programa detalhado (lista numerada de tópicos). Informe também curso, classificação, créditos, carga horária e objetivos.}
\begin{itemize}
    \item \textbf{Curso:} [curso];
    \item \textbf{Classificação:} [Teórica / Prática / Teórico-Prática];
    \item \textbf{Créditos:} [N];
    \item \textbf{Carga Horária:} [N horas];
    \item \textbf{Objetivos:} [objetivos da disciplina].
\end{itemize}

\subsubsection{Ementa}
[Ementa oficial da disciplina.]

\subsubsection{Programa Detalhado}
\begin{enumerate}
    \item [Tópico 1]
    \item [Tópico 2]
    \item [\ldots]
\end{enumerate}

\subsection{Bibliografia}
\instrucao{Liste a bibliografia básica e complementar adotada pelo professor, no formato ABNT (ou o padrão usado pelo plano de ensino).}
\textbf{Bibliografia Básica:}
\begin{itemize}
    \item [SOBRENOME, N. \textbf{Título}. Edição. Editora, ano.]
\end{itemize}

\textbf{Bibliografia Complementar:}
\begin{itemize}
    \item [SOBRENOME, N. \textbf{Título}. Edição. Editora, ano.]
\end{itemize}

\subsection{Formas de Avaliação}
\instrucao{Descreva a composição da nota e as datas das avaliações. Adapte a fórmula, a tabela e a lista de datas ao plano de ensino.}
A avaliação é composta pelos seguintes componentes:

\begin{equation*}
NF = [AV_1 + AV_2 + \ldots]
\end{equation*}

\vspace{0.3cm}
\noindent
\begin{tabular}{@{}p{2.7cm}p{4.1cm}p{3.6cm}p{1.8cm}@{}}
\hline
\textbf{Avaliação} & \textbf{Tipo} & \textbf{Valor} & \textbf{Peso} \\ \hline
{[AV1]} & [Prova / Trabalho] & [00,00 pontos] & [0,00] \\
{[AV2]} & [Prova / Trabalho] & [00,00 pontos] & [0,00] \\ \hline
\multicolumn{2}{l}{\textbf{Total}} & \textbf{[100,00 pontos]} & \\
\end{tabular}
\vspace{0.3cm}

\textbf{Datas das Avaliações:}
\begin{itemize}
    \item [Avaliação 1: DD/MM/AAAA]
    \item [Avaliação 2: DD/MM/AAAA]
\end{itemize}

% ========================== 3. ATIVIDADES ===================================
\section{Atividades do Estagiário} \label{sec:atividades}
\instrucao{Item 3 das Diretrizes: indique as atividades e tarefas que ficarão sob sua responsabilidade. Seja específico e realista; elas serão cobradas no Relatório Final. Exemplos comuns: elaboração e correção de trabalhos ou listas, atribuição de notas, esclarecimento de dúvidas, auxílio em aulas práticas e aplicação de provas, preparação de material didático, elaboração do relatório. Edite, remova ou acrescente itens.}
\begin{enumerate}
    \item [Elaboração de trabalhos / listas de exercícios];
    \item [Correção de trabalhos];
    \item [Atribuição de notas];
    \item [Esclarecimento de dúvidas dos alunos];
    \item [Auxílio ao professor nas aulas práticas e aplicações de prova];
    \item Elaboração e entrega do Relatório Final.
\end{enumerate}

% ============================ 4. CRONOGRAMA =================================
\section{Cronograma} \label{sec:cronograma}
\instrucao{Item 4 das Diretrizes: cronograma de aulas e avaliações do semestre, indicando sua participação ao longo do tempo. Inclua uma linha por aula ou evento (use ``---'' em feriados) e as datas de elaboração/correção dos trabalhos e de entrega do Relatório Final. Datas devem seguir o calendário acadêmico da UFMG.}

\begin{longtable}{|p{2.2cm}|p{7.6cm}|p{4.7cm}|}
\hline
\textbf{Data} & \textbf{Evento} & \textbf{Atividade de Estágio} \\
\hline
\endfirsthead
\hline
\textbf{Data} & \textbf{Evento} & \textbf{Atividade de Estágio} \\
\hline
\endhead
\hline
\endfoot

[DD/MM/AAAA] & [Aula teórica: tópico] & [Esclarecimento de dúvidas] \\
\hline
[DD/MM/AAAA] & [Feriado] & --- \\
\hline
[DD/MM/AAAA] & [Aula prática: tópico] & [Auxílio ao professor] \\
\hline
[DD/MM/AAAA] & [Prova escrita 1] & [Auxílio ao professor] \\
\hline
[DD/MM/AAAA] & [Aula teórica: tópico] & [Elaboração do Trabalho 1] \\
\hline
[DD/MM/AAAA] & [Monitoria / aula] & [Correção e lançamento de notas] \\
\hline
[DD/MM/AAAA] & [Encerramento das atividades] & [Entrega do Relatório Final] \\
\hline
\end{longtable}

% ======================== CONSIDERACOES FINAIS ==============================
\section{Considerações Finais}
\instrucao{1 parágrafo: declare que a proposta foi elaborada em conformidade com o plano de ensino oficial da disciplina no período indicado e resuma o objetivo formativo do estágio.}
Esta proposta de estágio foi elaborada em conformidade com o plano de ensino oficial da disciplina \NomeDisciplina{} (\CodigoDisciplina) oferecida no período \Semestre{} pela Universidade Federal de Minas Gerais. [Complemente com o objetivo formativo do estágio.]

\vspace{2.5cm}
\noindent
\begin{minipage}[t]{0.48\textwidth}
    \centering
    \rule{\linewidth}{0.5pt} \\
    \textbf{\NomeProfessor} \\
    \TituloProfessor \\
    Professor Supervisor da Disciplina
\end{minipage}
\hfill
\begin{minipage}[t]{0.48\textwidth}
    \centering
    \rule{\linewidth}{0.5pt} \\
    \textbf{\NomeEstagiario} \\
    \NivelEstagiario \\
    \EstagioNivel
\end{minipage}

\end{document}
